 \section{Related Work}\label{sec:related}

\subsection{Video Large Language Models}

Research on VideoLLMs has moved from simple feature alignment to more efficient token management and deeper reasoning in recent years. To balance spatial resolution and temporal coverage, Keye-VL-1.5~\citep{yang2025kwai} introduces a Slow-Fast video encoding strategy that dynamically allocates computation according to inter-frame similarity, enabling high-resolution modeling of key frames while preserving broader temporal coverage for relatively static content. 
Qwen3-VL~\citep{bai2025qwen3} enhances spatiotemporal modeling capabilities through interleaved-MRoPE, strengthens visual-language alignment with DeepStack, and introduces a text-based temporal alignment mechanism, thereby achieving more precise temporal localization in videos.
To reduce computational redundancy, InternVL3.5~\citep{wang2025internvl3} adopts Visual Resolution Routing and improves patch-level visual token compression without compromising performance, while LongVU~\citep{shen2024longvu} combines DINOv2~\citep{oquab2023dinov2} features with cross-modal queries to achieve adaptive spatiotemporal compression and mitigate memory bottlenecks in long-video processing. Recent studies also explore higher-level temporal reasoning and event modeling. VideoChat-A1~\citep{wang2026videochat} introduces the Chain-of-Shot reasoning paradigm to support stepwise reasoning over long videos, the VidEvent~\citep{liang2025videvent} dataset provides large-scale annotated data for modeling the evolution of dynamic events in video. Overall, these studies have advanced video modeling in terms of efficiency, reasoning, and temporal understanding, but most of them are still designed primarily for offline settings.

\subsection{Streaming Video Understanding}

The core of streaming video understanding lies in the shift from whole-video processing to online interactive processing~\citep{wang2025streambridge,xia2025streaming}. This shift requires the synchronization of perception, decision-making, and response, making streaming settings more challenging than conventional offline video understanding. Dispider~\citep{qian2025dispider} proposes a decoupled framework for perception and reaction, using a lightweight active streaming module to detect appropriate interaction points and avoid the latency introduced by autoregressive decoding during real-time monitoring. In proactive decision-making, MMDuet2~\citep{wang2025mmduet2} applies multi-round reinforcement learning and a PAUC reward mechanism to improve both response quality and timeliness without relying on precise timestamp annotations. MiniCPM-o-4.5~\citep{minicpmo45} extends omni-modal modeling to full-duplex live streaming by supporting simultaneous video and audio input together with concurrent text and speech output in an end-to-end framework. LiveCC~\citep{chen2025livecc} shows that low-cost ASR transcripts can serve as scalable supervision for streaming training, outperforming offline models of up to 72 billion parameters. StreamingVLM~\citep{xu2025streamingvlm} enables stable real-time understanding of infinite video streams through compact cache-based streaming inference and chunk-level supervised fine-tuning. VideoLLM-online~\citep{chen2024videollm} introduces a LIVE framework for streaming video dialogue, combining streaming-oriented training, dialogue data generation, and efficient real-time inference for continuous video understanding. These advances have been evaluated on benchmarks including OVO-Bench~\citep{niu2025ovo}, StreamingBench~\citep{lin2024streamingbench}, and OmniMMI~\citep{wang2025omnimmi}, which together form a more comprehensive benchmark suite for real-time interactive video understanding. Collectively, these studies highlight the importance of interaction-aware modeling to ensure both continuity and timeliness in streaming scenarios.
